\chapter [\CO{[KRI]}\\ \emph{Spezielle Patientengruppe:} Opfer von Kriminalität,
Misshandlung und Vernachlässigung] {\CC{[KRI]}\\ \emph{Spezielle
Patientengruppe:} Opfer von Kriminalität, Misshandlung und Vernachlässigung}

\Layout{2}{\TxQuerverweise}{%
\begin{inparaitem}
  \item Anzeigepflicht: \prettyref{topic:anzeigepflicht}
  \item Anzeigerecht: \prettyref{topic:anzeigerecht}
  \item Dokumentation: \ref{topic:dokumentation-rechtlich},
  \ref{topic:dokumentation-medizinisch}
\end{inparaitem}
}{}{}{}{}



\section{Handlungen gegen die sexuelle Selbstbestimmung und Integrität}

\subsection{\Z{Sexuelle Gewalt Vergewaltigung und geschlechtliche Nötigung}}

\AuthNotes{eher akut, besonderer Forkus Opferbetreuung}

Vergewaltigung
\ZFootnote{%
{\S}\,201. (1) Wer eine Person mit Gewalt, durch Entziehung der persönlichen
Freiheit oder durch Drohung mit gegenwärtiger Gefahr für Leib oder Leben
{\ldots} zur Vornahme oder Duldung des Beischlafes oder einer dem Beischlaf
gleichzusetzenden geschlechtlichen Handlung nötigt, ist mit Freiheitsstrafe {\ldots} zu bestrafen.

(2) {\ldots}
}

Geschlechtliche Nötigung
\ZFootnote{%
{\S}\,202. (1) Wer außer den Fällen des {\S}\,201 eine Person mit Gewalt oder
durch gefährliche Drohung zur Vornahme oder Duldung einer geschlechtlichen
Handlung nötigt, ist mit Freiheitsstrafe {\ldots} zu
bestrafen.

(2) {\ldots}
}


\Z{\TakeNotesLarge}

\subsection{\Z{Sexueller Missbrauch}}

\AuthNotes{eher chronisch, besonderer Fokus Erkennung und Dokumentation

umfasst grundsätzlich geschlechtliche Handlungen, inkl. Beischlaf und "`dem
Beischlaf gleichzusetzende Handlungen"', sowohl durch den Täter selbst
selbst, als auch die "`Verleitung"' zu Vornahme von Handlungen an oder durch
das Opfer selbst.}


Ebenso ist zu bestrafen, wer eine unmündige Person zur Vornahme oder Duldung des
Beischlafes \ldots dazu verleitet, eine dem Beischlaf gleichzusetzende
geschlechtliche Handlung an sich selbst vorzunehmen.

\begin{itemize}
  \item Sexueller Missbrauch einer wehrlosen oder psychisch beeinträchtigten Person
\ZFootnote{%
{\S}\,205. (1) Wer eine wehrlose Person oder eine Person, die wegen einer
Geisteskrankheit, wegen einer geistigen Behinderung, wegen einer tiefgreifenden
Bewusstseinsstörung oder wegen einer anderen schweren, einem dieser Zustände
gleichwertigen seelischen Störung \E{unfähig ist, die Bedeutung des Vorgangs
einzusehen} {\ldots}, \E{unter Ausnützung dieses
Zustands} dadurch missbraucht, dass er an ihr eine \E{geschlechtliche Handlung}
vornimmt oder von ihr an sich vornehmen lässt oder sie zu einer geschlechtlichen Handlung mit
einer anderen Person oder, um sich oder einen Dritten geschlechtlich zu erregen
oder zu befriedigen, dazu verleitet, eine geschlechtliche Handlung an sich selbst
vorzunehmen, ist mit Freiheitsstrafe von {\ldots} zu
bestrafen.

(2) {\ldots}
}
\item Schwerer sexueller Mißbrauch von Unmündigen
\ZFootnote{%
{\S}\,206. (1) Wer mit einer unmündigen Person den Beischlaf oder {\ldots} 
gleichzusetzende geschlechtliche Handlung unternimmt, ist mit Freiheitsstrafe {\ldots} zu bestrafen.

(2) Ebenso ist zu bestrafen, wer eine unmündige Person zur Vornahme oder Duldung
des Beischlafes oder {\ldots} gleichzusetzenden geschlechtlichen
Handlung mit einer anderen Person oder, um sich oder einen Dritten geschlechtlich
zu erregen oder zu befriedigen, dazu verleitet, eine {\ldots}
gleichzusetzende geschlechtliche Handlung an sich selbst vorzunehmen.

(3) {\ldots}

(4) {\ldots}
}
\item Sexueller Mißbrauch von Unmündigen
\ZFootnote{%
{\S}\,207. (1) Wer außer dem Fall des {\S}\,206 eine geschlechtliche Handlung an
einer unmündigen Person vornimmt oder von einer unmündigen Person an sich
vornehmen läßt, ist mit Freiheitsstrafe von sechs Monaten bis zu fünf Jahren zu
bestrafen.

(2) Ebenso ist zu bestrafen, wer eine unmündige Person zu einer geschlechtlichen
Handlung (Abs. 1) mit einer anderen Person oder, um sich oder einen Dritten
geschlechtlich zu erregen oder zu befriedigen, dazu verleitet, eine
geschlechtliche Handlung an sich selbst vorzunehmen.

(3) Hat die Tat eine schwere Körperverletzung ({\S}\,84 Abs. 1) zur Folge, so ist
der Täter mit Freiheitsstrafe von fünf bis zu fünfzehn Jahren, hat sie aber den
Tod der unmündigen Person zur Folge, mit Freiheitsstrafe von zehn bis zu zwanzig
Jahren oder mit lebenslanger Freiheitsstrafe zu bestrafen.

(4) Übersteigt das Alter des Täters das Alter der unmündigen Person nicht um mehr
als vier Jahre und ist keine der Folgen des Abs. 3 eingetreten, so ist der Täter
nach Abs. 1 und 2 nicht zu bestrafen, es sei denn, die unmündige Person hätte das
zwölfte Lebensjahr noch nicht vollendet. 
}
\item Sexueller Missbrauch von Jugendlichen
\ZFootnote{
{\S}\,207b. (1) Wer an einer Person, die das \E{16. Lebensjahr} noch nicht
vollendet hat und aus bestimmten Gründen noch nicht reif genug ist, {\ldots} unter Ausnützung dieser
mangelnden Reife sowie seiner altersbedingten Überlegenheit eine geschlechtliche
Handlung vornimmt, von einer solchen Person an sich vornehmen lässt oder eine
solche Person dazu verleitet, eine geschlechtliche Handlung an einem Dritten
vorzunehmen oder von einem Dritten an sich vornehmen zu lassen, ist mit
Freiheitsstrafe {\ldots} oder mit Geldstrafe {\ldots} zu
bestrafen.

(2) Wer an einer Person, die das 16. Lebensjahr noch nicht vollendet hat, unter
Ausnützung einer Zwangslage dieser Person eine geschlechtliche Handlung vornimmt,
von einer solchen Person an sich vornehmen lässt oder eine solche Person dazu
verleitet, eine geschlechtliche Handlung an einem Dritten vorzunehmen oder von
einem Dritten an sich vornehmen zu lassen, ist mit Freiheitsstrafe {\ldots} zu bestrafen.

(3) Wer eine Person, die das 18. Lebensjahr noch nicht vollendet hat, unmittelbar
durch ein Entgelt dazu verleitet, eine geschlechtliche Handlung an ihm oder einem
Dritten vorzunehmen oder von ihm oder einem Dritten an sich vornehmen zu lassen,
ist mit Freiheitsstrafe {\ldots} zu bestrafen. 
}
\item Sittliche Gefährdung von Personen unter sechzehn Jahren
\ZFootnote{
{\S}\,208. (1) Wer eine Handlung, die geeignet ist, die sittliche, seelische oder
gesundheitliche Entwicklung von Personen unter sechzehn Jahren zu gefährden, vor
einer unmündigen Person oder einer seiner Erziehung, Ausbildung oder Aufsicht
unterstehenden Person unter sechzehn Jahren vornimmt, um dadurch sich oder einen
Dritten geschlechtlich zu erregen oder zu befriedigen, ist mit Freiheitsstrafe
bis zu einem Jahr zu bestrafen, es sei denn, daß nach den Umständen des Falles
eine Gefährdung der unmündigen oder Person unter sechzehn Jahren ausgeschlossen
ist.

(2) Übersteigt das Alter des Täters im ersten Fall des Abs. 1 das Alter der
unmündigen Person nicht um mehr als vier Jahre, so ist der Täter nicht zu
bestrafen, es sei denn, die unmündige Person hätte das zwölfte Lebensjahr noch
nicht vollendet.
}
\item Mißbrauch eines Autoritätsverhältnisses
\ZFootnote{

{\S}\,212. (1) Wer
										
1. mit einer mit ihm in absteigender Linie verwandten minderjährigen Person, seinem
minderjährigen Wahlkind, Stiefkind oder Mündel oder

2. mit einer minderjährigen Person, die seiner Erziehung, Ausbildung oder Aufsicht
untersteht, unter Ausnützung seiner Stellung gegenüber dieser Person

eine geschlechtliche Handlung vornimmt oder von einer solchen Person an sich
vornehmen lässt oder, um sich oder einen Dritten geschlechtlich zu erregen oder
zu befriedigen, dazu verleitet, eine geschlechtliche Handlung an sich selbst
vorzunehmen, ist mit Freiheitsstrafe {\ldots} zu bestrafen.

(2) Ebenso ist zu bestrafen, wer
										
1. als Arzt, klinischer Psychologe, Gesundheitspsychologe, Psychotherapeut,
Angehöriger eines Gesundheits- und Krankenpflegeberufes oder Seelsorger mit einer
berufsmäßig betreuten Person,

2. als Angestellter einer Erziehungsanstalt oder sonst als in einer
Erziehungsanstalt Beschäftigter mit einer in der Anstalt betreuten Person oder

3. als Beamter mit einer Person, die seiner amtlichen Obhut anvertraut ist,
unter Ausnützung seiner Stellung dieser Person gegenüber eine geschlechtliche
Handlung vornimmt oder von einer solchen Person an sich vornehmen lässt oder, um
sich oder einen Dritten geschlechtlich zu erregen oder zu befriedigen, dazu
verleitet, eine geschlechtliche Handlung an sich selbst vorzunehmen.
}
\end{itemize}









\Z{\TakeNotesLarge}


\section{Gewaltdelikte}

\subsection{\Z{Häusliche Gewalt}}

\Z{\TakeNotesLarge}

\section{Missbrauch und Misshandlungen}

\subsection{Kindesmisshandlung}
\label{topic:kindesmisshandlung}

\Definition{Absichtliche psychische oder physische Gewalt gegen Kinder}


Kindesmisshandlung kommt in \Es{allen sozialen Schichten} vor!

\Layout{2}{Hinweiszeichen}
{
\begin{itemize}
  \item Langes Intervall zwischen  Verletzungszeitpunkt und Verständigung der
  Rettung.
  \item Anamnese und Verletzungsmuster  widersprechen einander!
  \item Mehrere Verletzungen verschiedenen Alters.
  \item Erkennbarer, verletzungsverursachender Gegenstand (Schuhabdruck,
  \ldots).
  \item Zeichen der Vernachlässigung.
  \item Verletzungen bzw. Schmerzen in der Genitalregion (Achtung: Es gibt auch andere
  Ursachen als sexueller Missbrauch!).
  \item "`Untypische"' Verletzungen.
  \item Verhaltensstörungen.
\end{itemize}
}{
\begin{itemize}
  \item Langes Intervall Verletzung -- Hilfeersuchen.
  \item Widerspruch Anamnese -- Verletzungsmuster.
  \item Verletzungen verschiedenen Alters.
  \item Erkennbarer, verletzungsverursachender Gegenstand .
  \item Vernachlässigung.
  \item Verletzungen / Schmerzen in Genitalregion (Achtung: auch andere
  Ursachen!).
  \item "`Untypische"' Verletzungen.
  \item Verhaltensstörungen.
\end{itemize}
}{}{}{}

\MassnahmenNG{Spezielle Maßnahmen}{Kindesmisshandlung}{}{
    Die Aufgaben des Rettungsdienstes bei Verdacht auf Kindesmisshandlung
besteht in:

\begin{enumerate}
  \item Daran denken!
  \item Umstände und Umfeld abklären
  \item Unstimmigkeiten und Auffälligkeiten auf eine neutrale Weise
  \Es{dokumentieren}
  \item Obiges verlässlich bei der Übergabe an die nachbetreuende Einrichtung
  weiterleiten. Evtl. Name des Gesprächpartners bei der  Übergabe dokumentieren.
  \item \Es{Patienten nicht im Stich lassen}: Gelindestes Mittel wählen,
  evtl. Vorwand für die Hospitalisierung suchen, aber wenn ein begründeter
  Verdacht auf eine Misshandlung besteht und die
  Erziehungsberechtigten unkooperativ sind notfalls auch die Exekutive beiziehen.
\end{enumerate}

\paragraph{Vorsicht!}

Jedenfalls zu \underline{unterlassen} sind:

\begin{itemize}
  \item Voreilige Verdachtsäußerungen. \Z{Nicht alles was stinkt ist ein
  Fisch!}
  \item Andeutungen gegenüber den Erziehungsberechtigten oder Dritten.
\end{itemize}

}



\section{Vernachlässigung}

\Layout{0}{Rechtliche Grundlagen}{%

Strafgesetzbuch: Quälen oder Vernachlässigen unmündiger, jüngerer
oder wehrloser Personen ({\S} 92)%
\ZFootnote{{\S}\,92 StGB: Quälen oder
Vernachlässigen unmündiger, jüngerer oder wehrloser Personen (Auszug):

\emph{(1) Wer einem anderen, der seiner Fürsorge oder Obhut untersteht und der
das achtzehnte Lebensjahr noch nicht vollendet hat oder wegen Gebrechlichkeit,
Krankheit oder einer geistigen Behinderung wehrlos ist, körperliche oder
seelische Qualen zufügt, ist mit Freiheitsstrafe bis zu drei Jahren zu
bestrafen.}

\emph{(2) Ebenso ist zu bestrafen, wer seine Verpflichtung zur Fürsorge oder
Obhut einem solchen Menschen gegenüber gröblich vernachlässigt und dadurch, wenn auch
nur fahrlässig, dessen Gesundheit oder dessen körperliche oder geistige
Entwicklung beträchtlich schädigt.}
}


}{
{\S} 92 StGB
}{}{}{}



